% Generated draft: requires mathematical and source review.
\documentclass[11pt]{article}
\usepackage[margin=25mm]{geometry}
\usepackage{fontspec}
\setmainfont{Latin Modern Roman}
\usepackage{amsmath,amssymb,mathtools,unicode-math}
\setmathfont{Latin Modern Math}
\usepackage{enumitem,needspace,xurl,hyperref}
\setlength{\parindent}{0pt}
\setlength{\parskip}{5pt}
\setlength{\emergencystretch}{3em}
\newcommand{\N}{\mathbb{N}}
\newcommand{\Z}{\mathbb{Z}}
\newcommand{\Q}{\mathbb{Q}}
\newcommand{\R}{\mathbb{R}}
\newcommand{\C}{\mathbb{C}}
\newcommand{\F}{\mathbb{F}}
\newcommand{\A}{\mathbb{A}}
\newcommand{\PP}{\mathbb{P}}
\newcommand{\E}{\mathbb{E}}
\newcommand{\eps}{\varepsilon}
\newcommand{\dd}{\,\mathrm{d}}
\begin{document}
\section*{UnsolvedMath ID: 30006493}
\textbf{Source problem code:} catalogue entry
\subsection{Definitions and mathematical statement}
The supplied record concerns solutions of the Yang--Baxter equation, welded pairs, structure skew braces, and deformations of solutions. It does not supply definitions of a welded pair or of the structure skew brace associated with a solution. In particular, it does not specify the underlying sets, the relevant operations, or the hypotheses imposed on solutions. The phrase ``the two families of solution deformations above'' refers to surrounding exposition that is absent from the supplied record.

Subject to recovering these definitions and the two deformation families from the cited report, the requested constructions are:
\begin{enumerate}
\item Determine an object associated with a welded pair that plays the same role as the structure skew brace associated with a solution.
\item Find a common framework embracing the two families of solution deformations referred to in the source.
\end{enumerate}
The properties characterizing ``the same role'' and the requirements for a ``common framework'' are not specified. Consequently, these questions preserve the supplied research targets but do not yet constitute a self-contained mathematical formulation. No domain, additional hypothesis, or universal property is imposed in this draft.
\subsection{Short English statement}
The problem asks for an analogue of the structure skew brace for welded pairs and for a framework encompassing two source-specific families of solution deformations. The supplied excerpt omits the definitions and deformation families needed to make either request precise. Its open status is supplied dated curation and requires independent verification.
\subsection{Sources}
\begin{itemize}
\item Local UnsolvedMath registry, supplied record with permanent ID \(30006493\), source problem code \texttt{catalogue entry}, title \emph{Structure Objects and Deformations of Welded Pairs}, in the collection \emph{Oberwolfach Reports Open Problems}. No individual proposer is identified; the supplied proposed year is 2025.
\item Supplied source citation: \emph{Mini-Workshop: The Yang--Baxter Equation and Representations of Braid Groups (2026)}, Oberwolfach Reports, pp.\ 2633--2672, DOI \texttt{10.4171/owr/2025/49}. \url{https://doi.org/10.4171/owr/2025/49}. The original statement includes a page marker for p.\ 2655. The citation's parenthetical year 2026 and the report designation 49/2025 are retained as supplied.
\item Supplied literature source: \emph{Oberwolfach Report 49/2025}, \url{https://doi.org/10.4171/owr/2025/49}. Literature checked at: \texttt{2026-08-22}. This is supplied dated curation, not independently verified: the record labels the questions open and reports that no subsequent general framework resolving either question was located. Its background explicitly calls for independent verification of time-sensitive status and cited claims.
\item The registry assigns statement status \texttt{corrected\_verified} and reports that the official report was checked, the question isolated from surrounding exposition, and its notation repaired. These are supplied source-verification claims, not verification performed for this draft. Oberwolfach Reports are described in the record as documenting mathematical workshops and their open-problem discussions.
\end{itemize}
\end{document}
